\chapter{Where to go next}
\label{app:reading}

Every chapter ended with a short list of books. This appendix steps back and suggests a path. It is built on a simple, widely shared conviction, put especially clearly by Kevin Zhou in his short guide \emph{The Theoretical Minimum}: you can get a deep understanding of physics by working through a small set of canonical textbooks, doing many of their problems yourself. The books below are recommended because they are self-contained, anticipate a learner's misconceptions, and above all contain good problems.

\section*{The undergraduate core}

If you work through these, doing most of the problems, you will know the core theoretical content of an undergraduate physics degree.
\begin{enumerate}[itemsep=4pt]
  \item \textbf{Any calculus-based introductory text} --- Halliday, Resnick, and Walker; Young and Freedman; or similar. They are about equally good; what matters is doing the problems.
  \item \textbf{John R.\ Taylor, \emph{Classical Mechanics}.} A smooth, thorough introduction to intermediate mechanics, starting from ordinary calculus. (Chapters~\ref{ch:newton}--\ref{ch:analytical}, \ref{ch:chaos}.)
  \item \textbf{Daniel V.\ Schroeder, \emph{An Introduction to Thermal Physics}.} Clean, focused, and clear. (Chapters~\ref{ch:thermo} and~\ref{ch:statmech}.)
  \item \textbf{David J.\ Griffiths, \emph{Introduction to Electrodynamics}.} The best undergraduate electromagnetism book by a wide margin, with an excellent crash course on vector calculus. (Chapter~\ref{ch:em}.)
  \item \textbf{Griffiths and Schroeter, \emph{Introduction to Quantum Mechanics}.} The most accessible introduction, built on elementary exercises plus a rich variety of harder problems. (Chapter~\ref{ch:qm}.)
\end{enumerate}
Alongside these, \textbf{Mary Boas, \emph{Mathematical Methods in the Physical Sciences}} (Chapter~\ref{ch:math}) supplies the mathematics as you need it.

\section*{Specializations at the undergraduate level}

Each of these is a standard first serious book in its field.
\begin{itemize}[itemsep=3pt]
  \item \textbf{Optics}: Hecht, \emph{Optics}.
  \item \textbf{Special relativity}: Taylor and Wheeler, \emph{Spacetime Physics}; or Morin, \emph{Special Relativity for the Enthusiastic Beginner}.
  \item \textbf{General relativity}: Hartle, \emph{Gravity}; or Schutz, \emph{A First Course in General Relativity}.
  \item \textbf{Fluids and continua}: Lautrup, \emph{Physics of Continuous Matter}.
  \item \textbf{Chaos}: Strogatz, \emph{Nonlinear Dynamics and Chaos}.
  \item \textbf{Computation}: Newman, \emph{Computational Physics}.
  \item \textbf{Condensed matter}: Simon, \emph{The Oxford Solid State Basics}.
  \item \textbf{Atomic physics}: Foot, \emph{Atomic Physics}; quantum optics: Gerry and Knight, \emph{Introductory Quantum Optics}.
  \item \textbf{Quantum information}: Nielsen and Chuang, \emph{Quantum Computation and Quantum Information}.
  \item \textbf{Nuclear and particle physics}: Krane, \emph{Introductory Nuclear Physics}; Griffiths, \emph{Introduction to Elementary Particles}.
  \item \textbf{Plasma physics}: Chen, \emph{Introduction to Plasma Physics and Controlled Fusion}.
  \item \textbf{Astrophysics}: Carroll and Ostlie, \emph{An Introduction to Modern Astrophysics}.
  \item \textbf{Cosmology}: Ryden, \emph{Introduction to Cosmology}.
  \item \textbf{Biophysics}: Nelson, \emph{Biological Physics}.
  \item \textbf{String theory}: Zwiebach, \emph{A First Course in String Theory}.
\end{itemize}

\section*{The graduate core}

A physics PhD typically begins by seeing the undergraduate core again at a higher level. Modern texts that improve on the long-standing classics include:
\begin{itemize}[itemsep=3pt]
  \item \textbf{Mechanics}: Helliwell and Sahakian, \emph{Modern Classical Mechanics}.
  \item \textbf{Electrodynamics}: Zangwill, \emph{Modern Electrodynamics}.
  \item \textbf{Statistical mechanics}: Kardar, \emph{Statistical Physics of Particles} and \emph{Statistical Physics of Fields}.
  \item \textbf{Quantum mechanics}: Robert Littlejohn's Berkeley lecture notes (free online), or Sakurai.
  \item \textbf{Quantum field theory}: Schwartz, \emph{Quantum Field Theory and the Standard Model}; or Peskin and Schroeder.
  \item \textbf{General relativity}: Carroll, \emph{Spacetime and Geometry}.
\end{itemize}

\section*{Problems, and more problems}

Physics is learned by solving problems. Beyond the textbooks above:
\begin{itemize}[itemsep=3pt]
  \item \textbf{David Morin}, \emph{Introduction to Classical Mechanics} and \emph{Problems and Solutions in Introductory Mechanics}: hundreds of beautiful problems with full solutions.
  \item \textbf{Kevin Zhou's physics olympiad handouts}: a free, carefully graded collection of challenging problems spanning mechanics, electromagnetism, thermodynamics, optics, and modern physics, with insightful commentary.
  \item \textbf{I.\ E.\ Irodov}, \emph{Problems in General Physics}: the classic Soviet problem book, terse and difficult.
  \item \textbf{Mahajan}, \emph{Street-Fighting Mathematics} and \emph{The Art of Insight in Science and Engineering}: problems in estimation and approximation --- the spherical cow, as a discipline.
\end{itemize}

\section*{For perspective}

\begin{itemize}[itemsep=3pt]
  \item \textbf{\emph{The Feynman Lectures on Physics}} (free online): the whole of physics, seen through the eyes of a master.
  \item \textbf{Feynman}, \emph{The Character of Physical Law} and \emph{QED}: short books on how physics works and what it means.
  \item \textbf{Steven Weinberg}, \emph{To Explain the World} and \emph{Dreams of a Final Theory}.
  \item \textbf{Free lecture notes}: David Tong's Cambridge notes on nearly every subject in theoretical physics; MIT OpenCourseWare; and, for what is happening right now, the physics section of the arXiv preprint server.
\end{itemize}

\section*{How to read a physics textbook}

Slowly, with a pencil. Before reading a derivation, try to guess the result; afterwards, close the book and reproduce it. Do the examples before reading the solutions. Do more problems than are assigned, and return to the ones you could not do. When you are stuck, estimate, find a limiting case, or draw a picture. And when you finish a book, you will find that the next one is easier than you feared. That is the point of the minimum: the foundation is set.
